\documentclass[10pt,a4paper]{amsart}
\usepackage[margin=19mm]{geometry}
\usepackage{amsmath,amssymb,mathtools}
\usepackage[scr=rsfs,cal=euler]{mathalfa}
\usepackage{xfrac}
\usepackage[unicode,hidelinks,bookmarks=false]{hyperref}
\newcommand{\ie}{i.e.\ }
\newcommand{\cf}{cf.\ }
\newcommand{\resp}{respectively\ }
\newcommand{\Subsection}{\subsection}
\providecommand{\nobreakdash}{\nobreak\hbox{-}}
\setlength{\emergencystretch}{2em}
\multlinegap0pt
\usepackage{bm}
\usepackage{bbm}
\usepackage{dsfont}
\usepackage{xspace}
\usepackage{cancel}
\usepackage{mathtools}
\usepackage{amsthm}
\makeatletter
\@ifundefined{theorem}{\newtheorem{theorem}{Theorem}}{}
\@ifundefined{lemma}{\newtheorem{lemma}[theorem]{Lemma}}{}
\makeatother
\title[Cops and Robbers, Clique Covers, and Induced Cycles]{Cops and Robbers, Clique Covers, and Induced Cycles}
\author{Alexander Clow, Imed Zaguia}
\date{Source: arXiv:2507.14321v1 (CC BY 4.0)}
\begin{document}
\maketitle

\noindent\emph{Source and licence.} Cops and Robbers, Clique Covers, and Induced Cycles. Version arXiv:2507.14321v1 (CC BY 4.0), distributed under the Creative Commons Attribution 4.0 International licence, as recorded in the arXiv licence metadata for this exact version and shown on the abstract page https://arxiv.org/abs/2507.14321v1. Full rights notice: licenses/arxiv-2507.14321-CC-BY-4.0.txt.\par\smallskip

\noindent\textit{Editorial scope.} Counted unit: the Lemma whose source label is \texttt{Lemma: Always an escape in H} in the article (source0.tex lines 306--318, source label \texttt{Lemma: Always an escape in H}), delivered together with the article's own complete original proof of it (source0.tex lines 320--382, one \texttt{proof} environment of 63 lines). No mathematics is written, repaired or re-derived: statement and proof are the article's own text line for line. Required context retained uncounted: none. Every definition, display and label of those lines is retained unchanged. Omitted from this package and not redistributed: 1--305, 319--319, 383--844. Nothing inside a retained excerpt is omitted or altered: every retained body is delivered exactly as the article's lines are written, so no omission receipt of any kind accompanies this package. Citations: the retained excerpts carry no citation command at all, so no bibliography entry of the article is transcribed and no reference is redistributed. Reference adaptation: none was needed and none was made; every reference inside the delivered unit resolves inside it, so no unresolved reference or citation is produced. Printed slips kept literal: no printed word, symbol, hypothesis or number of the counted statement or of its proof was altered. Layout and class: the article's own class is \texttt{amsart} with packages including amsfonts, graphicx, tabularx, array, color, comment, amsmath, amsthm, amssymb, fullpage, xcolor, tikz, listings, dsfont; the wrapper is the lane's pinned \texttt{amsart} editorial wrapper at 10pt on A4, which restates the theorem-like environments the retained text needs and adds the article's own macro definitions that the retained text needs (none), with extra wrapper packages (bm, bbm, dsfont, xspace, cancel, mathtools). The delivered numbering is the unit's own. Assets: no retained span contains a figure environment, a graphics-inclusion command, a PDF-page inclusion, a table or a TikZ picture; no image file is copied into this package, the package declares no asset dependency and no image file is redistributed with it. Licence: version arXiv:2507.14321v1 of this article is distributed under the Creative Commons Attribution 4.0 International licence, as recorded in the arXiv licence metadata for this exact version and shown on the abstract page https://arxiv.org/abs/2507.14321v1. A full rights notice is delivered at licenses/arxiv-2507.14321-CC-BY-4.0.txt. Characters: every retained excerpt is pure ASCII, so the delivered engine, XeLaTeX with its default Latin Modern text font, reports no missing character.
\begin{lemma}\label{Lemma: Always an escape in H}
    Let $k$ be a fixed constant, let $\ell\geq 2$ be fixed but arbitrary, and let $p = \frac{2k\log(\ell)}{\ell}$.
    If $v\in S_i$
    let $E_v$ be the event that there exists an integer $j\neq i$ and a set of vertices $u_1,\dots, u_{k-1} \in V(H(k,\ell,p))\setminus (S_j \cup \{v\})$
    such that 
    \[
    N(v) \cap S_j \subseteq \Big(\bigcup_{t=1}^{k-1} N(u_t) \Big).
    \]
    Then, the probability that no event $E_v$ occurs is at least 
    \[
    1 - \frac{k^{k+1}}{(k-1)!}\exp\Big(4 k^{3} \log(\ell)^2 \ell^{(\frac{1}{\ell}-1)}\Big)\ell^{-k}.
    \]
\end{lemma}

\begin{proof}
    Let $k$ be a fixed constant, let $\ell\geq 2$ be integer valued, let $p = \frac{2k\log(\ell)}{\ell}$,
    and 
    let $i$ and $v \in S_i \subseteq V(H(k,\ell,p))$ be fixed but arbitrary.
    For a fixed $j\neq i$, let $E_{v,j}$ be the event
    that there exists a set of vertices $u_1,\dots, u_{k-1} \in V(H(k,\ell,p))\setminus (S_j \cup \{v\})$
    such that 
    \[
    N(v) \cap S_j \subseteq \Big(\bigcup_{t=1}^{k-1} N(u_t) \Big).
    \]
    Note that by the symmetry of $H(k,\ell,p)$ the choice of $j$ does not effect the probability of the event $E_{v,j}$.
    Then, by the union bound $\mathbb{P}(E_v) \leq (k-1)\mathbb{P}(E_{v,j})$.
    So to bound the probability of $E_v$ it is sufficient to bound the probability of $E_{v,j}$ for a fixed $j$.

    Let $j$ be fixed.
    Let $X$ be the random variable that counts the cardinality of $N(v) \cap S_j$.
    Then $X$ has distribution $Bin(\ell,p)$.
    Hence,
    \begin{align*}
        \mathbb{P}(X = x) & = \binom{\ell}{x}p^x(1-p)^{\ell-x} \\
        & = (2k)^x\binom{\ell}{x}\log(\ell)^x \ell^{-x}\Big(1-\frac{2k\log(\ell)}{\ell} \Big)^{\ell-x} \\
        & \leq (2k)^x\frac{\ell^x}{x!}\log(\ell)^x \ell^{-x} \exp\Big(-2k\log(\ell) + \frac{2kx\log(\ell)}{\ell}\Big) \\
        & = \frac{(2k\log(\ell))^x}{x!}\ell^{(\frac{x}{\ell}-1)2k}.
    \end{align*}

    Let us now consider the probability of $E_{v,j}$ conditioned on the value of $X$.
    There are $\binom{(k-1)\ell-1}{k-1}$ choices of vertices $u_1,\dots, u_{k-1}$,
    and given a fixed choice of vertices $u_1,\dots, u_{k-1}$ 
    the probability of $N(v) \cap S_j \subseteq \Big(\bigcup_{t=1}^{k-1} N(u_t) \Big)$
    is at most $((k-1)p)^{X} = (\frac{2(k-1)k\log(\ell)}{\ell})^X$.

    By the law of total probability and the union bound 
    \[
     \mathbb{P}( E_{v,j}) \leq \binom{(k-1)\ell-1}{k-1} \sum_{x=0}^\ell \mathbb{P}\Big( E_{v,j} | X = x\Big)\mathbb{P}\Big(X = x\Big).
    \]
    So the product $\mathbb{P}( E_{v,j} | X = x)\mathbb{P}(X = x)$ is of interest.
    Let us consider this quantity; by the inequalities we have already shown
    \begin{align*}
        \mathbb{P}\Big( E_{v,j} | X = x\Big)\mathbb{P}\Big(X = x\Big) & \leq \Bigg(\frac{2(k-1)k\log(\ell)}{\ell} \Bigg)^x \Bigg(\frac{(2k\log(\ell))^x}{x!}\ell^{(\frac{x}{\ell}-1)2k} \Bigg) \\
        & < \ell^{-2k}\Bigg(\frac{2^{2x} k^{3x} \log(\ell)^{2x} \ell^{(\frac{1}{\ell}-1)x}}{x!} \Bigg).
    \end{align*}
    Plugging this into the earlier sum, and recalling the series definition of $e^z = \sum_{x=0}^\infty \frac{z^x}{x!}$,
    \begin{align*}
        \mathbb{P}( E_{v,j}) & \leq \binom{(k-1)\ell-1}{k-1} \sum_{x=0}^\ell \ell^{-2k}\Bigg(\frac{2^{2x} k^{3x} \log(\ell)^x \ell^{(\frac{1}{\ell}-1)x}}{x!} \Bigg) \\
        & < \ell^{-2k} \binom{(k-1)\ell-1}{k-1}  \sum_{x=0}^\infty \frac{\Big(4 k^{3} \log(\ell)^2 \ell^{(\frac{1}{\ell}-1)}\Big)^x}{x!} \\
        & = \ell^{-2k} \binom{(k-1)\ell-1}{k-1} \exp\Big(4 k^{3} \log(\ell)^2 \ell^{(\frac{1}{\ell}-1)}\Big) \\
        & < \ell^{-2k} \frac{(k\ell)^{k-1}}{(k-1)!} \exp\Big(4 k^{3} \log(\ell)^2 \ell^{(\frac{1}{\ell}-1)}\Big) \\
        & = \frac{k^{k-1}}{(k-1)!}\exp\Big(4 k^{3} \log(\ell)^2 \ell^{(\frac{1}{\ell}-1)}\Big)\ell^{-k-1}.
    \end{align*}
    Since $k$ is constant,
    observe the exponential term tends to $1$ as $\ell$ tends to infinity.
    Recalling $\mathbb{P}(E_v) \leq (k-1)\mathbb{P}(E_{v,j})$,
    we conclude that 
    \[
    \mathbb{P}(E_v) < \frac{k^{k}}{(k-1)!}\exp\Big(4 k^{3} \log(\ell)^2 \ell^{(\frac{1}{\ell}-1)}\Big)\ell^{-k-1}.
    \]
    Applying the union bound over all choice of $v$,
    the probability that there exists a vertex $v$ where event $E_v$ occurs is at most 
    \[
    \frac{k^{k+1}}{(k-1)!}\exp\Big(4 k^{3} \log(\ell)^2 \ell^{(\frac{1}{\ell}-1)}\Big)\ell^{-k}.
    \]
    This concludes the proof.
\end{proof}

\end{document}
